% Source PDF: source/普通高考/2015/2015全国2理(贵州,甘肃,青海,西藏,黑龙江,吉林,海南,宁夏,内蒙古,新疆,云南,广西,辽宁).pdf, page 1.
% The flowchart is vector line art; pdfimages -list found no complete embedded bitmap for it.
% Grid evidence: tmp/review_flowchart_2015_national_paper_2_science/grid/page-1-grid100.png
% (100px coarse + 50px fine) and q8-block-grid50.png (50px + 25px fine).
% Original comparison crop: tmp/review_flowchart_2015_national_paper_2_science/crop/q8_original_final.png,
% cropped from the ungridded page render.
% Elements: rounded start/end terminals, input/output trapezoids, a\ne b and a>b
% decision diamonds, subtraction process rectangles, directed connectors, branch labels,
% and the left return loop.
% Node -> directed-edge checklist: 开始 -> 输入 a,b -> (a != b); (a != b)-否 -> 输出 a -> 结束;
% (a != b)-是 -> (a > b); (a > b)-是 -> a=a-b; (a > b)-否 -> b=b-a;
% both subtraction boxes -> the left return loop -> the incoming stem of (a != b).
% solid segments: all node borders and directed connectors; dashed segments: none.
% labels: all node text is centered; 是/否 labels sit beside their corresponding decision exits;
% the return loop remains independent of the output branch and re-enters before (a != b).
\begin{tikzpicture}[
    x=1cm,
    y=1cm,
    >=Stealth,
    line width=0.45pt,
    line cap=round,
    line join=round,
    font=\normalsize,
    flow/.style={draw, anchor=center, align=center, inner xsep=4pt, inner ysep=2pt,
        execute at begin node={\setlength{\parindent}{0pt}}},
    branchlabel/.style={anchor=center, align=center, inner sep=0pt,
        execute at begin node={\setlength{\parindent}{0pt}}},
    terminal/.style={flow, rounded corners=0.18cm, minimum width=1.20cm, minimum height=0.58cm},
    process/.style={flow, rectangle, minimum height=0.66cm},
    decision/.style={flow, diamond, aspect=2.9, inner sep=0pt, minimum height=0.82cm},
    io/.style={flow, trapezium, trapezium left angle=72, trapezium right angle=108,
        minimum height=0.66cm},
]
    % Main path and output branch.
    \node[terminal] (start) at (0,7.80) {开始};
    \node[io, minimum width=2.15cm] (input) at (0,6.70) {输入 $a,b$};
    \node[decision, minimum width=3.15cm] (neq) at (0,5.60) {$a\ne b$};
    \node[io, minimum width=1.85cm] (output) at (2.75,4.15) {输出 $a$};
    \node[terminal] (finish) at (2.75,3.10) {结束};

    % Subtraction branch decisions and process boxes.
    \node[decision, minimum width=2.85cm] (gt) at (-2.00,4.20) {$a>b$};
    \node[process, minimum width=2.05cm] (suba) at (-3.35,3.10) {$a=a-b$};
    \node[process, minimum width=2.05cm] (subb) at (-0.65,3.10) {$b=b-a$};

    \draw[->] (start.south) -- (input.north);
    \draw[->] (input.south) -- (neq.north);
    \draw[->] (neq.east) -- (2.75,5.60) -- (2.75,4.48) -- (output.north);
    \draw[->] (output.south) -- (finish.north);

    % 是 branch from a != b to the comparison decision.
    \draw[->] (neq.west) -- (-2.00,5.60) -- (gt.north);
    % The a>b branches enter their own subtraction boxes.
    \draw[->] (gt.west) -- (-3.35,4.20) -- (suba.north);
    \draw[->] (gt.east) -- (-0.65,4.20) -- (subb.north);

    % Both subtraction results merge into the independent return loop.
    \coordinate (merge) at (-2.00,2.55);
    \draw (suba.south) -- (-3.35,2.55) -- (merge);
    \draw (subb.south) -- (-0.65,2.55) -- (merge);
    \draw[->] (merge) -- (-2.00,1.85) -- (-4.90,1.85) -- (-4.90,6.20) -- (0,6.20);

    \node[branchlabel] at (-1.35,5.95) {是};
    \node[branchlabel] at (1.35,5.95) {否};
    \node[branchlabel] at (-3.05,4.50) {是};
    \node[branchlabel] at (-1.00,4.50) {否};
\end{tikzpicture}
